\input{../tex-inputs/latex-leseansicht-vorspann}

\section[Emil Marriot an Arthur Schnitzler, 1. 5. 1902]{L04333 Emil Marriot an Arthur Schnitzler, 1. 5. 1902}
\nopagebreak\mylabel{L04333v}
\rehead{ }\normalsize\beginnumbering\briefempfaengerindex{Schnitzler, Arthur@\textsc{Schnitzler, Arthur}!zzzMarriot, Emil@\emph{von Emil Marriot}!1902-05-011@{1. 5. 1902}|(be}
\toendnotes[C]{\smallbreak\pagebreak[2]}
\correspDesc{Versand  durch Emil Marriot am 1. 5. 1902 in Wien
\newline{}Erhalt  durch Arthur Schnitzler im Zeitraum [1. 5. 1902 – 4. 5. 1902?] in Wien}\toendnotes[C]{\smallbreak}
\Standort{DLA, A:Schnitzler, HS.1985.1.4008.}
\physDesc{Brief, 1 Blatt, 2 Seiten, 707 Zeichen
\newline{}Handschrift: schwarze Tinte, lateinische Kurrent
\newline{}Schnitzler: mit Bleistift Vermerk: »\textsc{Schüttelgasse 31, I}\oindex{XXXX Ortsangabe fehlt|pw}« }\toendnotes[C]{\smallbreak}
\pstart
           {\pb}Wien\oindex{Wien@\textbf{Wien}, \emph{Verwaltungsgebiet}|pw}, 1. Mai 1902.\pend
           \vspace{0.5em}
\pstart
           Sehr geehrter Herr, ich freue mich, zu hören, dass wir \label{K_L04333-1v}\edtext{uns »eigentlich« bereits kennen}{\lemma{\textnormal{\emph{uns … kennen}}}\Cendnote{\textnormal{Sie hatten sich am
                     20. 3. 1881 kennengelernt, vgl. A. S.: \emph{Tagebuch}, 22. 3. 1881.}}}\label{K_L04333-1}, wenn
               ich’s auch vergessen hatte. Gott! muss ich damals bei Zuckers\pwindex{Zucker, Karl @\textsc{Zucker, Karl}|pw}\pwindex{Zucker, Joseph 1821/1822 – 10.\,9.\,1890 Wien@\textsc{Zucker, Joseph} (1821/1822 – 10.\,9.\,1890 Wien), \emph{Advokat}|pw}\pwindex{Lorm, Jenny @\textsc{Lorm, Jenny}, \emph{Schauspielerin}|pw}\pwindex{\textcolor{red}{\textsuperscript{XXXX indx1}}|pw} unausstehlich gewesen sein!
               Wollen sie es trotzdem riskiren und mir die Freude machen, zu mir zu kommen?
               Vielleicht an einem Nachmittag zum Thee, so um 5 Uhr? Ich würde mich sehr freuen. Der
               Tag wäre mir für die nächste Zeit ganz gleich. Nur würde {\pb}ich Sie bitten, mich vorher wissen zu lassen, an welchem Tag ich Sie erwarten darf,
               um Sie nicht zu verfehlen.\pend
           
\pstart
           Denken Sie: ich kenne »Freiwild\pwindex{Schnitzler, Arthur 15. 5. 1862 Wien – 21. 10. 1931 ebd.@\textsc{Schnitzler, Arthur} (15. 5. 1862 Wien – 21. 10. 1931 ebd.), \emph{Schriftsteller, Mediziner}!Freiwild. Schauspiel in 3 Akten@\strich\emph{Freiwild. Schauspiel in 3 Akten}|pw}« noch nicht. Aber jetzt
               will ich’s mir ansehen. Das »\label{K_L04333-2v}\edtext{deutsche Theater\orgindex{Deutsches Theater Berlin@Deutsches Theater Berlin|pw}« bringt es ja diesmal}{\lemma{\textnormal{\emph{deutsche … diesmal}}}\Cendnote{\textnormal{Das \emph{Berliner deutsche Theater}\orgindex{Deutsches Theater Berlin@Deutsches Theater Berlin|pwk} führte \emph{Freiwild}\pwindex{Schnitzler, Arthur 15. 5. 1862 Wien – 21. 10. 1931 ebd.@\textsc{Schnitzler, Arthur} (15. 5. 1862 Wien – 21. 10. 1931 ebd.), \emph{Schriftsteller, Mediziner}!Freiwild. Schauspiel in 3 Akten@\strich\emph{Freiwild. Schauspiel in 3 Akten}|pwk} auf
                  seinem Gastspiel in Wien\oindex{Wien@\textbf{Wien}, \emph{Verwaltungsgebiet}|pwk}, das am 6. 5. 1902 beginnen sollte, nicht auf.}}}\label{K_L04333-2}, glaube ich.\pend
           
\pstart
           Mit guten Grüssen{\\[\baselineskip]}Ihre ganz ergebene{\\[\baselineskip]}\spacefill\mbox{E. Marriot.}\pend
           \leftskip=0em{}\selectlanguage{ngerman}\endnumbering\briefempfaengerindex{Schnitzler, Arthur@\textsc{Schnitzler, Arthur}!zzzMarriot, Emil@\emph{von Emil Marriot}!1902-05-011@{1. 5. 1902}|)be}\mylabel{L04333h}
\begin{anhang}
\end{anhang}\newcommand{\dateiname}{L04333}\newcommand{\titel}{Emil Marriot an Arthur Schnitzler, 1. 5. 1902}\newcommand{\editorInnen}{Herausgegeben von Jahnke, SelmaMüller, Martin Anton}\input{../tex-inputs/latex-leseansicht-abspann}
